\documentclass[aspectratio=169,10pt,t]{beamer}
\usepackage{ac1}
\aulainfo{14}{Exercícios e Revisão}{Revisão geral dos capítulos 1 a 4 (preparação para a Prova 1)}

\begin{document}
\capa

\begin{frame}{Mapa da matéria da Prova 1}
  \centering\small
  \begin{tabular}{p{2.6cm}p{6cm}l}
    \toprule
    \textbf{Capítulo} & \textbf{Tópicos principais} & \textbf{Aulas}\\
    \midrule
    1. Computadores & bases 2/8/10/16, complemento a 2, IEEE 754, ASCII, von Neumann, compilado \emph{vs.} interpretado & 2, 3\\
    2. Bases & algoritmo, fluxograma, pseudocódigo, identificadores, tipos, constantes & 4, 5, 6\\
    3. Estruturada & \texttt{if}/\texttt{elif}/\texttt{else}, \texttt{while}, \texttt{for}, laços encaixados & 7, 8, 9\\
    4. Subalgoritmos & funções, parâmetros, escopo, módulos, \texttt{try}/\texttt{except} & 10 a 13\\
    \bottomrule
  \end{tabular}
\end{frame}

\begin{frame}{Cap.~1: o que você precisa saber fazer}
  \begin{itemize}
    \item Converter entre \textbf{decimal, binário, octal e hexadecimal} (à mão)
    \item Representar inteiros com sinal em \textbf{complemento a 2} e somar em binário
    \item Explicar por que \(0{,}1+0{,}2\neq0{,}3\) (IEEE 754)
    \item Diferença entre \textbf{kilo} e \textbf{kibi}; calcular tempos de download e capacidades
    \item Desenhar o modelo de \textbf{von Neumann} e explicar CPU, memória, E/S
    \item Diferenciar linguagens \textbf{compiladas}, \textbf{interpretadas} e \textbf{semi-interpretadas}
  \end{itemize}
  \begin{exampleblock}{Exercícios do site (cap.~1)}
    1.3 a 1.9 (conversões), 1.10, 1.11 (unidades), 1.12 a 1.15 (complemento a 2), 1.16 a 1.18 (reais)
  \end{exampleblock}
\end{frame}

\begin{frame}{Mini-teste: Capítulo 1}
  \begin{enumerate}
    \item Converta \(237_{10}\) para binário e hexadecimal.
    \item Qual a faixa de inteiros representáveis com 12 bits em complemento a 2?
    \item Escreva \(-18\) em complemento a 2 com 8 bits.
    \item Quantos MiB cabem em um arquivo de 5\,000\,000 bytes? (\(1\text{ MiB}=2^{20}\) bytes)
    \item Por que \texttt{0.1 + 0.2 == 0.3} dá \texttt{False}?
  \end{enumerate}
  \vspace{1ex}
  \small\textbf{Respostas:} (1) \(11101101_2=\text{ED}_{16}\); (2) \(-2048\) a \(2047\); (3) \texttt{1110 1110}; (4) \(\approx 4{,}77\) MiB; (5) 0,1 e 0,2 não têm representação binária exata.
\end{frame}

\begin{frame}{Cap.~2: o que você precisa saber fazer}
  \begin{itemize}
    \item Escrever um \textbf{algoritmo} em pseudocódigo e desenhar o \textbf{fluxograma}
    \item Decidir se um nome é um \textbf{identificador} válido
    \item Reconhecer tipos: \texttt{int}, \texttt{float}, \texttt{str}, \texttt{bool}, \texttt{complex}
    \item Usar \texttt{input} com conversão de tipo e \texttt{print} formatado
    \item Calcular expressões respeitando a \textbf{precedência}: \texttt{**}, \texttt{*}, \texttt{/}, \texttt{//}, \texttt{\%}
  \end{itemize}
  \begin{block}{Mini-teste: o que cada linha imprime?}
    \texttt{print(7 // 2, 7 \% 2, 7 / 2, 2 ** 3 ** 2, "5" + "3", int("5") + 3)}
  \end{block}
  \small \textbf{Resposta:} \texttt{3 1 3.5 512 53 8}
\end{frame}

\begin{frame}[fragile]{Cap.~3: leia o código e diga a saída}
\begin{lstlisting}
soma = 0
for i in range(1, 6):
    if i % 2 == 0:
        continue
    soma += i
print(soma)

n = 20
while n > 0:
    n -= 7
print(n)
\end{lstlisting}
  \begin{itemize}
    \item Primeiro laço: soma dos ímpares de 1 a 5 \(\to\) \textbf{9}
    \item Segundo: \(20\to13\to6\to-1\) \(\to\) \textbf{-1}
    \item Faça o ``teste de mesa'': uma tabela com os valores das variáveis a cada volta
  \end{itemize}
\end{frame}

\begin{frame}[fragile]{Cap.~3: escreva o código}
  \begin{enumerate}
    \item Leia um inteiro \(n>0\) e mostre a soma dos \textbf{múltiplos de 3} entre 1 e \(n\).
    \item Leia notas até o usuário digitar \(-1\); mostre a média e quantas notas foram \(\ge 7\).
    \item Escreva a condição ``o ano \texttt{a} é bissexto''.
    \item Imprima o padrão (laços encaixados):
  \end{enumerate}
\begin{lstlisting}[style=saida]
1
1 2
1 2 3
1 2 3 4
\end{lstlisting}
  \small Soluções comentadas de problemas semelhantes: exercícios 3.7 a 3.15 do site.
\end{frame}

\begin{frame}[fragile]{Cap.~4: leia o código (escopo e parâmetros)}
\begin{lstlisting}
x = 10

def f(a, b=2):
    x = a * b
    return x

def g():
    print(x)

print(f(3), f(3, 4), f(b=5, a=1))
g()
print(x)
\end{lstlisting}
  \begin{itemize}
    \item \texttt{f(3)} \(\to\) 6; \texttt{f(3, 4)} \(\to\) 12; \texttt{f(b=5, a=1)} \(\to\) 5
    \item \texttt{g()} imprime \textbf{10}: o \texttt{x} de \texttt{f} é local
    \item Último \texttt{print}: \textbf{10}
  \end{itemize}
\end{frame}

\begin{frame}[fragile]{Cap.~4: escreva o código}
  \begin{enumerate}
    \item Função \texttt{eh\_primo(n)} que devolve \texttt{True} ou \texttt{False}.
    \item Função \texttt{estatisticas(a, b, c)} que devolve o \textbf{menor}, o \textbf{maior} e a \textbf{média} dos três números.
    \item Função \texttt{le\_inteiro\_positivo(msg)} que só aceita inteiros \(>0\) (use \texttt{try}/\texttt{except}).
  \end{enumerate}
  \vspace{1ex}
\begin{lstlisting}
def eh_primo(n):
    if n < 2:
        return False
    for d in range(2, int(n ** 0.5) + 1):
        if n % d == 0:
            return False
    return True
\end{lstlisting}
  \small Teste: \texttt{eh\_primo(1)}, \texttt{eh\_primo(2)}, \texttt{eh\_primo(97)}, \texttt{eh\_primo(91)}.
\end{frame}

\begin{frame}{Erros que mais custam nota}
  \begin{columns}[T]
    \column{0.5\textwidth}
    \begin{itemize}
      \item Esquecer de converter o resultado de \texttt{input}
      \item \texttt{=} no lugar de \texttt{==}
      \item \texttt{range(n)} versus \texttt{range(1, n+1)}
      \item Acumulador que não é reiniciado
      \item Dividir sem proteger o zero
    \end{itemize}
    \column{0.5\textwidth}
    \begin{itemize}
      \item \texttt{print} no lugar de \texttt{return}
      \item Usar variável local fora da função
      \item Recuo (indentação) inconsistente
      \item Comparar \texttt{float} com \texttt{==}
      \item Não testar os limites das faixas
    \end{itemize}
  \end{columns}
\end{frame}

\begin{frame}{Orientações para a Prova 1 (Aula 15)}
  \begin{itemize}
    \item \textbf{Conteúdo:} capítulos 1 a 4 do livro
    \item Parte teórica à mão (conversões, complemento a 2, teste de mesa)
    \item Parte de programação: escrever funções e programas completos
    \item Estude refazendo os exercícios do site \emph{sem} olhar a solução
    \item Use os exercícios adicionais para treinar mais:
          \par\href{\siteurl extras/}{\texttt{\siteurl extras/}}
  \end{itemize}
  \begin{alertblock}{Prova 1}
    Aula 15: sem consulta. Chegue com tempo e releia o enunciado duas vezes.
  \end{alertblock}
\end{frame}

\begin{frame}{Lista de exercícios para revisão (site do livro)}
  \small
  \begin{tabular}{@{}p{2.4cm}p{9.3cm}@{}}
    \toprule
    \textbf{Capítulo 1} & 1.3 a 1.9, 1.10, 1.11, 1.12 a 1.18\\
    \textbf{Capítulo 2} & 2.1 a 2.8, 2.9 a 2.18\\
    \textbf{Capítulo 3} & 3.1 a 3.15\\
    \textbf{Capítulo 4} & 4.1 a 4.17 (pelo menos 4.7, 4.8, 4.11, 4.12, 4.13, 4.16)\\
    \bottomrule
  \end{tabular}
  \vspace{1ex}

  \href{\siteurl exercicios/}{\texttt{\siteurl exercicios/}}
\end{frame}

\end{document}
